\documentclass[11pt]{article}
\usepackage[hmargin=0.6in,vmargin=0.5in]{geometry}
\usepackage{amsmath}
\usepackage{amssymb}
\usepackage{amsthm}
\usepackage{microtype}
\usepackage[hidelinks]{hyperref}
\usepackage{enumitem}
\setlength{\parskip}{0pt}
\setlist{nosep,leftmargin=1.5em}

\newcommand{\R}{\mathbb{R}}
\newcommand{\N}{\mathbb{N}}
\newcommand{\supp}{\operatorname{supp}}
\newcommand{\Diff}{\operatorname{Diff}}
\newcommand{\Mo}{\mathring{M}}
\newcommand{\Mb}{\breve{M}}
\DeclareMathOperator{\ind}{ind}

\title{\large A note on the topology-coincidence statements in\\ \emph{Topology of the space of conormal distributions}}
\author{\normalsize Kunal Tyagi\thanks{\texttt{hello.jay.tyagi@gmail.com}. \textbf{Use of AI.} The mathematics in this note (the derivations, the computations, and the text) was produced by Claude (Anthropic) working under the author's direction; the author set the objectives, adjudicated the technical and editorial decisions, and is responsible for the content. Claude is a tool and is not an author.}}
\date{\vspace{-1ex}\normalsize September 6, 2026}

\begin{document}
\maketitle
\vspace{-2em}

\noindent\textbf{1. Purpose.}
While using the results of \cite{ALKL24} (cited by its published numbering and pages; arXiv:2304.00798v2 and v3 carry the same numbering, and every proposition, corollary, remark and equation cited here from \S\S3--4 has the same number and wording in all three arXiv versions), the author found that several statements asserting that two topologies coincide on a whole step of a filtration are false, already in the simplest cases, and that some later proofs depend on them. No erratum could be found (\S5). This note is offered to the authors as a courtesy: it lists the statements (\S2), gives explicit witnesses with the computation visible (\S3), and records what survives and how the affected proofs can be re-routed (\S4). The main package of the paper (acyclicity and retractivity of $I$, $A$, $\dot A$, $\mathcal K$, $J$, $K$ for compact $M$, and the exact sequences) stands once this is done.

\medskip\noindent\textbf{2. The statements.}
The following are false as stated: Prop.~3.2 (p.~16); Cor.~3.4, both assertions (p.~17); the acyclicity and bounded-retractivity clauses of Cor.~3.6 (p.~18) for every non-compact $U$ and every $l\ge1$, together with the part of the sentence on p.~18 that extends Prop.~3.2 and Cor.~3.4 to symbols on a vector bundle, and its extension of Cor.~3.6 when the base manifold is non-compact (the extensions of Prop.~3.3 and Cor.~3.5 are unaffected, and the extension of Cor.~3.6 over a compact base is true, \S4); Cor.~4.5 (p.~23), for every compact $(M,L)$ with $L\ne\emptyset$ and $\operatorname{codim}L\ge1$; Prop.~6.12 and Cor.~6.14, both assertions (p.~39); Cor.~6.21 (p.~40); Cor.~7.13, both assertions (p.~56). Cor.~6.27 (p.~41) and Cor.~7.22 (p.~58) are true, but their printed proofs pass through Cor.~6.21, hence through Cor.~4.5. The witnesses (a)--(e) below are constant in the base variable up to a fixed cut-off, so they also refute every parametrized or bundle version; (f) is inherently non-local.

\medskip\noindent\textbf{3. Witnesses.}
Notation as in \cite{ALKL24}: $S^m=S^m(U\times\R^l)$ with the seminorms $\|\cdot\|_{K,\alpha,\beta,m}$ of (3.1), $\|\cdot\|_{Q,C^k}$ of (3.4), $\|\cdot\|'_{K,\alpha,\beta,m}$ of (3.5); every (3.4) and (3.5) seminorm is continuous on $S^m$. Fix $\theta\in C_c^\infty(\R^l)$ with $\theta(0)=1$ and $\supp\theta\subset B(0,r)$.

\emph{(a) Prop.~3.2 and the first assertion of Cor.~3.4.} Let $g_N(x,\xi)=\theta(\xi-Ne_1)$. Then $g_N\in S^{-\infty}\subset S^m$ for every $m$. If $Q\subset K\times\bar B(0,\rho)$, then $\|g_N\|_{Q,C^k}=0$ for $N>\rho+r$, and every seminorm (3.5) of $g_N$ vanishes because $g_N$ has compact support in $\xi$. So $g_N\to0$ in the topology of (3.4) and (3.5), and $g_N\to0$ in $C^\infty(U\times\R^l)$. But for $|\beta|>m$ with $\partial^\beta\theta\not\equiv0$ (such $\beta$ exist in every order, a compactly supported nonzero function being no polynomial),
\[
\|g_N\|_{K,0,\beta,m}=\sup_{|\eta|\le r}|\partial^\beta\theta(\eta)|\,(1+|\eta+Ne_1|)^{|\beta|-m}\ \ge\ (1+N-r)^{|\beta|-m}\sup|\partial^\beta\theta|\longrightarrow\infty .
\]
Hence $(g_N)$ is unbounded in $S^m$: the topology of (3.4) and (3.5) is strictly coarser than that of $S^m$, for every $m\in\R$, every $U$ and every $l\ge1$. Taking $|\beta|>m'$: $g_N\in S^m$, $g_N\to0$ in $C^\infty$, $g_N\not\to0$ in $S^{m'}$.

\emph{(b) Where the proof of Prop.~3.2 fails.} Let $S'^m$ denote $S^m$ with the topology of (3.4) and (3.5). Fix $R>2r$, $c_i=(1+iR)^{m+1}$, and $b_j=\sum_{1\le i\le j}c_i\,\theta(\xi-iRe_1)$ (disjoint supports). For $i<j$, $b_j-b_i$ has compact $\xi$-support inside $\{|\xi|\ge(i+1)R-r\}$, so its (3.5) seminorms vanish and its (3.4) seminorms over a given $Q$ vanish for $i$ large: $(b_j)$ is Cauchy in $S'^m$. Its $C^\infty$-limit is the locally finite sum $b=\sum_ic_i\theta(\xi-iRe_1)$, and $|b(x,iRe_1)|\,(1+iR)^{-m}=1+iR\to\infty$, so $b\notin S^m$ and $(b_j)$ has no limit in $S'^m$: $S'^m$ is not complete, contrary to ``Hence $S'^m(U\times\R^l)$ is complete'' (p.~17). Concretely, for the class $a$ of $(b_j)$ in the completion one has $\|a\|'_{K,0,0,m}=\lim_j\|b_j\|'_{K,0,0,m}=0$ while $\|\varphi(a)\|'_{K,0,0,m}=\|b\|'_{K,0,0,m}=\infty$; the identity ``$\|\varphi(a)\|'_{K,\alpha,\beta,m}=\|a\|'_{K,\alpha,\beta,m}<\infty$'' at the top of p.~17 is where the argument breaks.

\emph{(c) Second assertion of Cor.~3.4.} Let $c_j=e^j\theta(\xi-je_1)$, so $c_j\in S^m$ for all $m$ and $c_j\to0$ in $C^\infty$. Fix $x_0\in U$, $K=\{x_0\}$, and for $k\in\N$ put $\varepsilon_k=\tfrac12\inf_{j\ge1}e^j(1+j)^{-k}>0$ and $W_k=\{a\in S^k:\|a\|_{K,0,0,k}<\varepsilon_k\}$. The absolutely convex hull $W$ of $\bigcup_kW_k$ meets every step $S^k$ in a $0$-neighborhood, so it is a $0$-neighborhood of $S^\infty=\ind_kS^k$. If $c_j=\sum_i\lambda_iw_i$ with $\sum_i|\lambda_i|\le1$ and $w_i\in W_{k_i}$, evaluation at $(x_0,je_1)$ gives $e^j=|c_j(x_0,je_1)|<\sum_i|\lambda_i|\,\varepsilon_{k_i}(1+j)^{k_i}\le\tfrac12e^j$, a contradiction. So $c_j\notin W$ for all $j$, and $c_j\not\to0$ in $S^\infty$: the topologies of $S^\infty$ and $C^\infty$ differ on every $S^m$. (No regularity of $S^\infty$ is used.)

\emph{(d) Cor.~4.5.} Let $M$ be compact, $L\ne\emptyset$, $\operatorname{codim}L=n'\ge1$, $\dim L=n''$, $m<m'<m''$, and $\bar m,\bar m',\bar m''$ as in (4.9). Take a chart $(U_1,x=(x',x''))$ adapted to $L$ and meeting $L$, $x\colon U_1\to U'\times U''$ (p.~19), $g\in C_c^\infty(U'')$, $g\not\equiv0$, and the partition of unity $\{h,f_i\}$ of (4.10) with $f_1=1$ and $h=f_i=0$ ($i\ge2$) on a neighborhood $N$ of the compact set $x^{-1}(\{0\}\times\supp g)\subset L$ (replace any $\{h,f_i\}$ by $\{(1-\rho)h,\ \rho+(1-\rho)f_1,\ (1-\rho)f_i\ (i\ge2)\}$, $\rho\in C_c^\infty(U_1)$, $\rho=1$ on $N$; by the closed graph theorem the Fr\'echet topology of $I^m(M,L)$ is independent of these choices anyway). Let $\psi\in C_c^\infty(\{|x'|<1\})$, $\psi\not\equiv0$; $N_0\in\N_0$ with $2N_0\ge\bar m''$; $\psi_*=\Delta^{N_0}\psi$, so $\hat\psi_*(\eta)=(-|\eta|^2)^{N_0}\hat\psi(\eta)$ is Schwartz and $\partial^\beta\hat\psi_*$ vanishes to order $2N_0-|\beta|$ at $0$. Put
\[
u_j(x',x'')=j^{\,n'+\bar m'}\,g(x'')\,\psi_*(jx').
\]
For $j$ large, $\supp u_j\subset x^{-1}(\{|x'|\le1/j\}\times\supp g)\subset N$, so $u_j\in C_c^\infty(U_1)\subset C^\infty(M)\subset I^{m}(M,L)$, $hu_j=0$ and $f_iu_j=0$ ($i\ge2$), and by (4.8) the symbol of $f_1u_j=u_j$ is $a_j(x'',\xi)=j^{\bar m'}g(x'')\hat\psi_*(\xi/j)$. Substituting $\xi=j\eta$,
\[
\|a_j\|_{K,\alpha,\beta,\bar m_1}=j^{\,\bar m'-\bar m_1}\,\sup_K|\partial^\alpha g|\;\Phi_{\beta,\bar m_1}(j),\qquad
\Phi_{\beta,\bar m_1}(j)=\sup_\eta|\partial^\beta\hat\psi_*(\eta)|\,(1/j+|\eta|)^{|\beta|-\bar m_1}.
\]
For $\bar m_1\le2N_0$ one has $\Phi_{\beta,\bar m_1}(j)\le C_\beta$ (use the Schwartz decay if $|\beta|\ge\bar m_1$, the vanishing at $0$ if $|\beta|<\bar m_1$), while $\Phi_{0,\bar m'}(j)\ge|\hat\psi_*(\eta_0)|\min\{(1+|\eta_0|)^{-\bar m'},|\eta_0|^{-\bar m'}\}=c>0$ for any $\eta_0\ne0$ with $\hat\psi_*(\eta_0)\ne0$. Hence $\|a_j\|_{K,\alpha,\beta,\bar m''}\le Cj^{\,\bar m'-\bar m''}\to0$, i.e.\ $u_j\to0$ in $I^{m''}(M,L)$, whereas, for any compact $K$ meeting $\{g\ne0\}$, $\|a_j\|_{K,0,0,\bar m'}\ge c\sup_K|g|>0$, i.e.\ $u_j\not\to0$ in $I^{m'}(M,L)$. The same $u_j$, with $\supp\psi\subset(\tfrac12,1)$ in the collar chart of $(\Mb,\partial M)$, lies in $C_c^\infty(\Mo)\subset\dot A^{m}(M)=I^m_M(\Mb,\partial M)$ (6.47) and refutes Cor.~6.21.

\emph{(e) Prop.~6.12, Cor.~6.14, Cor.~7.13.} Let $M$ be compact with boundary, $(x,y)$ a collar chart, $\chi\in C_c^\infty((\tfrac12,2))$ with $\chi(1)=1$, $g\in C_c^\infty$ in $y$ with $g(y_0)=1$, $m'\in\R$, and $u_j=j^{-m'}\chi(jx)g(y)$. Then $u_j\in C_c^\infty(\Mo)\subset A^{m_1}(M)$ for every $m_1$; $u_j\to0$ in $C^\infty(\Mo)$ (the supports leave every compact set); every seminorm (6.42) of $u_j$ vanishes for $j$ large, and every seminorm (6.43) vanishes for every $j$ ($u_j=0$ on $\{x<1/(2j)\}$). But the seminorm (6.41) with $P=1$ (the topology of $A^m(M)$ is the projective one over all $P\in\Diff_b(M)$, p.~38) gives $\sup_{\Mo}x^{-m'}|u_j|\ge(1/j)^{-m'}j^{-m'}\chi(1)g(y_0)=1$. So (with $m'=m$) $u_j\to0$ in the topology of (6.42) and (6.43) on $A^m(M)$ but not in $A^m(M)$, and (with $m'<m$) $u_j\to0$ in $C^\infty(\Mo)$ but not in $A^{m'}(M)$. For the second assertion of Cor.~6.14 take $v_j=e^j\chi(jx)g(y)$ and the $0$-neighborhood of $A(M)=\bigcup_kA^{-k}(M)$ built as in (c) from $W_k=\{u\in A^{-k}(M):\sup x^k|u|<\tfrac12\inf_{i\ge1}e^ii^{-k}\}$; evaluation at $(1/j,y_0)$ shows $v_j\notin W$, so $v_j\not\to0$ in $A(M)$ while $v_j\to0$ in $C^\infty(\Mo)$. Since $J^m(M,L)$ is given by (7.27) with $\Diff(M,L)$ and $\boldsymbol x^mL^\infty(M)$, $\boldsymbol x$ the extension of $|x|$ to $M$ whose lift is a boundary defining function of the manifold with boundary $\boldsymbol M$ obtained by cutting $M$ along $L$ (p.~56), the same $u_j$ and $v_j$, placed in $\{\tfrac12<jx<2\}\subset M\setminus L$, refute both assertions of Cor.~7.13. (\cite{ALKL24m} restates these second assertions at \S2.5.10, p.~38, and \S2.6.7, p.~53.)

\emph{(f) Cor.~3.6 for non-compact $U$.} Let $U\subset\R^n$ be open and nonempty, $n\ge1$, and let $l\ge1$. Choose $x_j\in U$ leaving every compact subset of $U$, disjoint closed balls $\bar B(x_j,r_j)\subset U$ forming a locally finite family, $\chi_j\in C_c^\infty(B(x_j,r_j))$ with $\chi_j(x_j)=1$, and $b_j(x,\xi)=\chi_j(x)(1+|\xi|^2)^{j/2}\in S^j\setminus S^m$ ($m<j$). Then $\{b_j\}$ is bounded in $S^\infty$: given an absolutely convex $0$-neighborhood $W$, pick in each $W\cap S^k$ a basic neighborhood $W_k=\{a\in S^k:\|a\|_{K_k,\alpha,\beta,k}<\varepsilon_k,\ (\alpha,\beta)\in F_k\}$ ($F_k$ finite). For $j$ so large that $\supp\chi_j\cap K_1=\emptyset$, write $b_j=b_j\rho_R+b_j(1-\rho_R)$ with $\rho_R(\xi)=\rho(\xi/R)$, $\rho\in C_c^\infty$, $\rho=1$ near $0$. The first term lies in $S^{-\infty}\subset S^1$ and vanishes on $K_1\times\R^l$, so $2b_j\rho_R\in W_1$; the second lies in $S^{j+1}$, vanishes for $|\xi|\le R$, and its $S^{j+1}$ seminorms are $O(1/R)$ by the Leibniz rule, so $2b_j(1-\rho_R)\in W_{j+1}$ for $R$ large; by absolute convexity $b_j\in W$. The finitely many remaining $b_j$ are absorbed by $W$. Thus $\{b_j\}$ is bounded in $S^\infty(U\times\R^l)$ but contained in no step: the spectrum is not regular, hence not boundedly retractive and (by \cite[Cor.~6.5]{Wen03} as quoted on p.~5) not acyclic; the printed derivations of the other clauses of Cor.~3.6 pass through these properties. The same applies to the extension on p.~18 for bundles over a non-compact manifold, and to \cite[\S2.1.8, p.~15]{ALKL24m}, which restates Cors.~3.4--3.6 for arbitrary open $U\subset\R^n$. (Sect.~4.3.3, p.~23, excepts acyclicity for non-compact $M$; the symbol statements do not.)

\medskip\noindent\textbf{4. What survives, and how.}
For a compact $K\subset U$ let $S^m_K$ be the closed subspace of $S^m(U\times\R^l)$ of symbols vanishing for $x\notin K$, with the seminorms $N_m(a;\gamma)=\|a\|_{K,\alpha,\beta,m}$, $\gamma=(\alpha,\beta)$. For $m<m'<m''$ and every $\gamma$ there are a finite set $\Gamma$ of multi-indices and constants $C,R_0>0$ such that
\begin{equation}\label{interp}
N_{m'}(a;\gamma)\ \le\ C\Big[R^{-(m'-m)/2}\,N_m(a;0)+R^{\,m''-m'}\max_{\gamma'\in\Gamma}N_{m''}(a;\gamma')\Big]\qquad(a\in S^{m''}_K,\ R\ge R_0),
\end{equation}
proved by Taylor expansion (an inequality of Landau--Kolmogorov type) on boxes of side $\sim(1+|\xi|)^{-c}$ in $x$ and $(1+|\xi|)^{1-c}$ in $\xi$ in the region $1+|\xi|>R$; compact base support makes $N_m(a;0)$ a global supremum, which is exactly what fails in (f). By \eqref{interp}, the $S^{m'}_K$- and $S^{m''}_K$-topologies coincide on the $0$-neighborhood $\{N_m(\cdot\,;0)<1\}$ of $S^m_K$: this is the criterion of \cite[Thm.~6.1]{Wen03} quoted on p.~4, so $\bigcup_mS^m_K$ is acyclic. (The inclusions $S^m_K\to S^{m'}_K$ also map bounded sets to relatively compact sets, by Arzel\`a--Ascoli and the tail bound $(1+|\xi|)^{m-m'}$, which is what the Montel clauses need; they are not compact operators, so the remark on p.~5 about compact spectra does not apply as it stands.) The same inequality holds for the collar models of $A^m(M)$, with weight $x^{-m}$, $x=e^{-\varrho}$, and $b$-derivatives. Now $I^m(M,L)$ carries by definition the initial topology of (4.10), whose symbols vanish outside the compact base projections of the $f_i$; $A^m(M)$ carries the initial topology of $u\mapsto((\lambda_iu)_i,\mu u)$ into collar models and $C^\infty(\Mo)$, for a partition of unity $\{\lambda_i,\mu\}$ subordinate to collar charts and $\Mo$ (by (6.33) and the local form of $\Diff_b(M)$); $\dot A^m(M)$, $\mathcal K^m(M)$, $K^m(M,L)$ are closed subspaces of $I^m(\Mb,\partial M)$, $\dot A^m(M)$, $I^m(M,L)$ ((6.47), pp.~41, 58); and $J^m(M,L)\cong A^m(\boldsymbol M)$ by (7.26). The criterion passes to initial topologies of a fixed map, to finite products of spectra each satisfying it, and to subspaces. Hence, for compact $M$, the LF-spaces $I(M,L)$, $A(M)$, $\dot A(M)$, $\mathcal K(M)$, $J(M,L)$, $K(M,L)$, whose steps are Fr\'echet spaces, are acyclic, hence complete, regular and boundedly, compactly and sequentially retractive by the results of \cite{Wen03} quoted on pp.~4--5. This gives valid proofs of the acyclicity clauses of Cor.~3.6 (compact base, e.g.\ $S^\infty(\R^l)$), 4.7, 6.16, 6.22, 6.28, 7.15 and 7.23, and then of their Montel clauses.

Several printed arguments use a coincidence statement for more than the word ``acyclic''. The semi-Montel step of the printed proof of Cor.~3.6 (p.~18) applies Cor.~3.4 to a bounded subset $B$ of some $S^m$; what it needs is the bounded-set statement, which is true: for $m'>m$, $S^{m'}$- and $C^\infty$-convergence agree on bounded subsets of $S^m(U\times\R^l)$, so, with $m'>m$ in place of $m$ and bounded retractivity in hand, the step goes through. In the proof of Claim~6.46 (pp.~48--49), the sentence ``By Corollary 6.14, there is some 0-neighborhood $V\subset A(M)$ such that $V\cap A^m(M)=W\cap A^m(M)$'' can be replaced as follows: the absolutely convex hull of $A$ is bounded in some $A^m(M)$, on which the topologies of $A(M)$, $A^{m'}(M)$ and $A^m(M)$ coincide by bounded retractivity, so there is $V$ with $V\cap A\subset W$; this suffices, with $E_{m'}$ from Prop.~6.29 (whose domain is $A^{m'}(M)\supset A^m(M)$, so $m'<m$ rather than $m'>m$ is the index needed). Prop.~8.8 (p.~64) is repaired identically with Cor.~7.31; Cor.~7.13 is not needed there. The density Cors.~6.15 and 7.14, whose printed proofs are modeled on Cor.~3.5's, follow instead from Prop.~6.10 and its $J$-analogue, or from Cors.~6.39 and 7.17 (Remarks~6.17 and 6.41 and their analogues on p.~57). Cor.~3.5 follows directly: for $a\in S^m$ and $m<m'$, $\|a-\rho(\xi/R)a\|_{K,\alpha,\beta,m'}\le CR^{m-m'}$ with the $\rho$ of (f), by the Leibniz rule, and a cut-off in $x$ finishes. The conclusion of Remark~3.8 (p.~18) likewise follows from (c) directly, without the sequential retractivity of $S^\infty(U\times\R^l)$, which is what its printed argument uses. Finally, Cor.~6.27 and Cor.~7.22 are true: by Prop.~7.26 every element of $K(M,L)$ is a finite sum of Dirac layers $\partial_x^k\delta_L\otimes v_k$, whose local symbols are polynomials in $\xi$ with the $v_k$ as coefficients, so it lies in $K^m(M,L)$ iff $v_k=0$ for $k>\bar m$; on polynomial symbols of bounded degree $k_0$ every $S^{\bar m_1}$-topology with $\bar m_1\ge k_0$ is the $C^\infty$-topology of the coefficients; Cor.~6.27 is the case $(\Mb,\partial M)$ via (6.49). In the memoir \cite{ALKL24m}, the uses of compact retractivity of $I(\mathcal F)=I(M,M^0;\Lambda\mathcal F)$ and $J(\mathcal F)$ (\S5.2.1, p.~119; \S\S5.5.3--5.5.4, p.~122) concern a closed $M$ and are covered by the above.

\medskip\noindent\textbf{5. Closing.}
As of September 6, 2026 the author found no erratum: arXiv lists versions v1 to v3 only, the Crossref record of the DOI carries no correction, zbMATH 7901419 carries no review or corrigendum, and Semantic Scholar lists no citing work (all four re-checked on that date); Unpaywall listed no correction DOI on September 3, 2026. Full derivations of (a)--(f), of \eqref{interp}, and of \S4 are available on request. The author would be glad to learn whether the authors agree, and thanks them for a paper the author has found very useful.

\vspace{-1em}
\begin{thebibliography}{9}\setlength{\itemsep}{0pt}\footnotesize
\bibitem{ALKL24} J.~A. \'Alvarez L\'opez, Y.~A. Kordyukov, E. Leichtnam, Topology of the space of conormal distributions, \emph{J. Pseudo-Differ. Oper. Appl.} 15 (2024), article 47, 68 pp., \url{https://doi.org/10.1007/s11868-024-00617-y}; arXiv:2304.00798 (v3, 1 June 2024).
\bibitem{ALKL24m} J.~A. \'Alvarez L\'opez, Y.~A. Kordyukov, E. Leichtnam, A trace formula for foliated flows, arXiv:2402.06671 (v2, 13 February 2024; the pages cited are the same in v1).
\bibitem{Wen03} J. Wengenroth, \emph{Derived Functors in Functional Analysis}, Lecture Notes in Mathematics, vol.~1810, Springer, Berlin (2003).
\end{thebibliography}
\end{document}
